\section{Proof Certificates}
\label{sec:method}
We propose proof certificates for LTL verification: \emph{transition safety certificates} and \emph{transition persistence certificates}. To verify that a system $\mathfrak{S}$ satisfies an LTL formula $\phi$, we transform the verification problem into proving that accepting transitions in the NBA for $\neg\phi$ occur finitely often. By the definition of NBA acceptance, a word is accepted if and only if accepting transitions occur infinitely often along some run. If we prove all accepting transitions occur only finitely often for any trajectory in $\mathfrak{T}(\mathfrak{S})$, then no word in $\mathfrak{L}(\mathfrak{S})$ can be accepted by the NBA for $\neg\phi$, thereby establishing $\mathfrak{S} \models_{\mathfrak{L}} \phi$. Both certificates ensure that accepting transitions in the NBA (representing $\neg\phi$) occur only finitely often. The safety certificate proves accepting transitions never occur by blocking reachability to their start states. The persistence certificate allows such states to be reachable but proves transitions occur finitely often via a ranking function.

\subsection{Transition Safety Certificate}

We prove accepting transitions never occur by showing their start states are unreachable in the product system.

\begin{definition}[transition safety certificate]
  \label{transition-safety-certificate-definition}
  Given a system $\mathfrak{S} = (\mathcal{X}, \mathcal{X}_0, f)$, a labelling function $\mathfrak{L}:\mathcal{X} \rightarrow \Sigma$ and an NBA $\mathcal{A}=(\Sigma, Q, q_0, \delta, Acc)$. Construct the product system $\mathfrak{S} \times \mathcal{A} = (\mathcal{Y}, \mathcal{Y}_0, F)$, where $\mathcal{Y} = \{(x, q) \,|\, x \in \mathcal{X}, q \in Q \}$, $\mathcal{Y}_0 = \{ (x_0, q_0) \,|\, x_0 \in \mathcal{X}_0 \}$, and $F(x,q) = \{(x^{\prime}, q') \,|\, x^{\prime} = f(x) \land (q, \mathfrak{L}(x), q') \in \delta\}$. For $q_k \in Acc^{start}$, define the unsafe set $\mathcal{Y}_u = \{(x, q_k) \,|\, x \in \mathcal{X}\}$.

  A bounded function $B_{k}(x, q): \mathcal{X} \times Q \rightarrow \mathbb{R}$ is a \emph{transition safety certificate} with respect to $q_k \in Acc^{start}$ if for all $(x,q) \in \mathcal{Y}$, $(x_0, q_0) \in \mathcal{Y}_0$, $(x^{\prime}, q^{\prime}) \in F(x, q)$, we have:
  \begin{align}
    &B_{k}(x_0, q_0) \geq 0, \label{btsc1} \\
    &B_{k}(x, q_k) < 0, \label{btsc2} \\
    &B_{k}(x, q) \geq 0 \Rightarrow B_{k}(x^{\prime}, q^{\prime}) \geq 0. \label{btsc3}
  \end{align}
\end{definition}

The certificate $B_k$ forms a barrier in the product system: non-negative at initial states (\ref{btsc1}), negative at all $(x, q_k)$ (\ref{btsc2}), and invariant under transitions (\ref{btsc3}). This guarantees $(x, q_k)$ is unreachable, preventing all accepting transitions from $q_k$.

\begin{lemma}[transition safety]
  \label{transition-safety-lemma}
  If a transition safety certificate $B_k$ with respect to $q_k \in Acc^{start}$ can be found, then all runs of words in $\mathfrak{L}(\mathfrak{S})$ never visit state $q_k$, thereby preventing all accepting transitions starting from $q_k$.
\end{lemma}

\subsection{Transition Persistence Certificate}

We prove accepting transitions occur finitely often using a ranking function that decreases with each occurrence.

\begin{definition}[transition persistence certificate]
  \label{transition-persistence-certificate-definition}
  Given a system $\mathfrak{S} = (\mathcal{X}, \mathcal{X}_0, f)$, a labelling function $\mathfrak{L}:\mathcal{X} \rightarrow \Sigma$ and an NBA $\mathcal{A}=(\Sigma, Q, q_0, \delta, Acc)$. For $(k, l) \in Acc^{pair}$, define the indicator function $I_{k,l}(x) = 1$ if $(q_k, \mathfrak{L}(x), q_l) \in Acc^{k,l}$, and $0$ otherwise.

  A bounded function $B_{k, l}(x): \mathcal{X} \rightarrow \mathbb{R}$ is a \emph{transition persistence certificate} with respect to $Acc^{k,l}$ if there exists $\epsilon > 0$ such that for all $x_0 \in \mathcal{X}_0$, $x \in \mathcal{X}$, $x^{\prime} = f(x)$, we have:
  \begin{align}
    &B_{k, l}(x_0) \geq 0, \label{btpc1} \\
    &B_{k, l}(x) \geq 0 \Rightarrow B_{k, l}(x^{\prime}) \geq 0, \label{btpc2} \\
    &B_{k, l}(x) \geq B_{k, l}(x^{\prime}) + I_{k,l}(x)\epsilon. \label{btpc3}
  \end{align}
\end{definition}

The certificate $B_{k,l}$ is non-negative initially (\ref{btpc1}), remains non-negative along trajectories (\ref{btpc2}), and decreases by $\epsilon$ whenever the state's label matches an accepting transition from $q_k$ to $q_l$ (\ref{btpc3}). By (\ref{btpc1}) and (\ref{btpc2}), $B_{k,l}$ remains non-negative along any trajectory starting from $x_0 \in \mathcal{X}_0$. By (\ref{btpc3}), $B_{k,l}$ decreases by at least $\epsilon$ each time $I_{k,l}(x) = 1$. Since accepting transitions from $q_k$ to $q_l$ can only occur when $(q_k, \mathfrak{L}(x), q_l) \in Acc^{k,l}$, the number of such transitions is bounded by $B_{k,l}(x_0)/\epsilon$, ensuring finite occurrence.

\begin{lemma}[transition persistence]
  \label{transition-persistence-lemma}
  If a transition persistence certificate $B_{k,l}$ can be found, then all runs of words in $\mathfrak{L}(\mathfrak{S})$ make accepting transitions from $q_k$ to $q_l$ happen finitely often.
\end{lemma}

\subsection{Main Verification Theorem}

We now present the main theorem for LTL verification. For each accepting transition start state in the NBA, we use either a transition safety certificate or transition persistence certificates to prove accepting transitions from that state occur finitely often.

\begin{theorem}[LTL verification]
  \label{ltl-verification-theorem}
   Given a system $\mathfrak{S} = (\mathcal{X}, \mathcal{X}_0, f)$, a labelling function $\mathfrak{L}: \mathcal{X} \rightarrow \Sigma$ and an NBA $\mathcal{A} = (\Sigma, Q, q_0, \delta, Acc)$ representing the negation of an LTL formula $\phi$. For each $q_k \in Acc^{start}$, if we can find either:
  \begin{enumerate}
      \item A transition safety certificate $B_k$ with respect to $q_k$, or
      \item Transition persistence certificates $B_{k,l}$ for all $(k, l) \in Acc^{pair}$ where $q_k$ is the start state,
  \end{enumerate}
  then it can be guaranteed that $\mathfrak{S} \models_{\mathfrak{L}} \phi$.
\end{theorem}

\begin{proof}
  Given an accepting transition start state $q_k \in Acc^{start}$:
  \begin{enumerate}
      \item If a transition safety certificate $B_k$ is found, by Lemma~(\ref{transition-safety-lemma}), all runs never visit state $q_k$, thereby preventing all accepting transitions starting from $q_k$.

      \item For each $(k, l) \in Acc^{pair}$ where the start state is $q_k$, if a transition persistence certificate $B_{k,l}$ is found, by Lemma~(\ref{transition-persistence-lemma}), all runs make accepting transitions from $q_k$ to $q_l$ happen finitely often.
  \end{enumerate}

  This ensures that all accepting transitions in $Acc$ occur only finitely often. By the definition of NBA acceptance, a word $w$ is accepted if and only if there exists a run where accepting transitions occur infinitely often, i.e., $\mathit{Inf}(r) \cap Acc \neq \emptyset$. Since all accepting transitions occur only finitely often, for any word $w \in \mathfrak{L}(\mathfrak{S})$ and any run $r$ on $w$, we have $\mathit{Inf}(r) \cap Acc = \emptyset$. Therefore, no word in $\mathfrak{L}(\mathfrak{S})$ is accepted by automaton $\mathcal{A}$, implying $\mathfrak{L}(\mathfrak{S}) \cap \mathfrak{L}(\mathcal{A}) = \emptyset$. Since $\mathfrak{L}(\mathcal{A}) = \mathfrak{L}(\neg \phi)$, we have $\mathfrak{L}(\mathfrak{S}) \subseteq \mathfrak{L}(\phi)$, and consequently $\mathfrak{S} \models_{\mathfrak{L}} \phi$.
\end{proof}

The two certificates are complementary: safety certificates block reachability to prove transitions never occur, while persistence certificates allow reachability but prove finite occurrence via ranking functions. This combination reduces conservatism while maintaining smaller search spaces than closure certificates. To illustrate their complementarity, consider verifying $\mathbf{G} \neg a$ with the NBA for $\mathbf{F} a$ (Figure \ref{pic-tbnba1}), where accepting transitions $(q_0, \emptyset, q_0)$ and $(q_0, \{a\}, q_0)$ form a self-loop. For the persistence certificate $B_{0,0}$, conditions (\ref{btpc1}-\ref{btpc3}) require $B_{0,0}(x_0) \geq 0$, invariance under transitions, and $B_{0,0}(x) \geq B_{0,0}(x^{\prime}) + \epsilon$. In this self-loop, the accepting transition occurs at every step, forcing infinite descent—contradicting boundedness. The persistence certificate fails for infinitely-occurring transitions. However, a safety certificate can prove $q_0$ is unreachable in the product system. Conversely, if $q_0$ is reachable but accepting transitions occur finitely often, the persistence certificate succeeds. This demonstrates their complementary strengths.

\subsection{Compared with Other Methods}
In the context of this discussion, we assume that the NBA for an LTL specification $\phi$ has already been constructed. The translation from an LTL formula to an NBA involves an exponential complexity with respect to the size of the formula \cite{vardi2005automata}. The complexity of constructing the complement of an NBA is known to be EXPTIME \cite{complementnba}.

\textbf{The State-Triplet Approach.} The state-triplet approach \cite{wongpiromsarn2015automata} uses barrier certificates (\ref{SCDEF}) to separate safe regions $\mathcal{X}_{safe}$ from unsafe regions $\mathcal{X}_{u}$. It blocks all simple paths from $q_0$ to accepting states in the NBA for $\neg\phi$. Finding certificates for all such paths ensures $\mathfrak{L}(\mathfrak{S})$ never visits accepting states, verifying $\mathfrak{S} \models_{\mathfrak{L}} \phi$. Our transition safety certificate shares this reachability-blocking principle but operates on the product system $\mathfrak{S} \times \mathcal{A}$, eliminating the need to enumerate simple paths. The state-triplet method requires one certificate per triplet to block each path individually, while we synthesize one certificate per accepting transition start state. Our persistence certificate further allows visiting accepting states while proving transitions occur finitely often. Combining both certificates reduces conservatism: beyond blocking reachability, we verify systems where runs reach accepting states yet satisfy the specification.

\textbf{Closure Certificates.} Closure certificates (\ref{LTLCC}) are defined on $\mathcal{X} \times Q \times \mathcal{X} \times Q$, creating high-dimensional search spaces with heavy computational burdens. Our certificates operate on smaller spaces: persistence certificates on $\mathcal{X}$, safety certificates on $\mathcal{X} \times Q$. We require $|Acc^{pair}|$ persistence certificates and $|Acc^{start}|$ safety certificates, yielding a search space of at most $\mathcal{X} \times Q \times Q$—smaller than $\mathcal{X} \times Q \times \mathcal{X} \times Q$. Case studies show that for benchmarks requiring closure certificates \cite{murali2024closure}, our certificates synthesize in seconds with significant speedup.
